% !TeX root = ../paper.tex
\hypertarget{motivating-diagnostics-why-vlas-struggle-with-diverse-instructions-and-reliable-control}{%
\section{Motivational Diagnostics}\label{motivating-diagnostics-why-vlas-struggle-with-diverse-instructions-and-reliable-control}}

The behavioral interventions reveal a gap between original-task competence and responses to local
commands, visual prompts, and corrections. To determine what can be used to bridge this gap, we ask
whether the relevant semantic and visual information remains accessible after action training. If it
is absent, improving the mapping from existing representations to actions alone would be
insufficient; the relevant representations would also need to be learned. We train target-point
probes on frozen representations under paired language and visual interventions. Target assignment
rates reach 83.80--98.82\%, above the 50\% two-choice reference, and predicted targets shift
consistently with the interventions. These findings show that target-relevant information remains
readable despite execution failures. Readability does not establish that the policy uses this
information, but provides a basis for learning more explicit correspondences between current
instructions, visual targets, and local actions. This motivates the supervision and training design
below, rather than relying solely on additional trajectory fitting. Full protocols, metric
definitions, and results are provided in Appendices~\ref{appendix:behavior-diagnostics}
and~\ref{appendix:representation-diagnostics}.

